\documentclass{exam}

\usepackage[margin=1in]{geometry}

\usepackage{comment}
\usepackage{graphicx}
\usepackage{listings}

\pagestyle{headandfoot}
\firstpageheader{CSCI 221}{}{Advanced Assembly Assignment}
\runningheader{CSCI 221}{}{Advanced Assembly Assignment}


\begin{document}
\title{Advanced Assembly Assignment (Evaluation)}
\author{CSCI 221: Computer Science Fundamentals 2}
\date{Spring 2026}
\maketitle


This assignment is an opportunity to test your understanding of advanced assembly. 
Feel free to collaborate with others and use resources, but all code and writeups must be your own. 
To submit your assignment, please submit your code and writeup. 
Mark which sections of the code are for which questions by using different files or comments in the file. 
For this assignment, your code does not need to check for invalid input. 
For this assignment, you are allowed to use pseudoinstructions. 

\textbf{Due Date:}
Friday, April 9th at 1:00 pm. 

For this assignment, you will be writing code in MIPS. 
Unfortunately, most machines you interact with will use a different ISA than MIPS, which means that they will be unable to run MIPS code. 
In order to test and debug your code, you will need to make use of a MIPS simulator. 
Many MIPS simulators are available online. 
If you have a preferred simulator, feel free to use it. 
If not, we recommend the QtSPIM MIPS simulator. 

\begin{questions}

\question[12]
\textbf{Decaying Sum.}
\begin{parts}
\part[6]
Rewrite the following code to be tail recursive. 
You may choose to write additional helper functions. 
\begin{verbatim}
uint32_t decaying_sum(uint16_t* values, uint16_t length, uint16_t decay){
    if (length <= 0){
        return 0;
    }
    uint32_t rest = decaying_sum(&values[1], length-1, decay);
    uint32_t decayed = rest / decay;
    return values[0] + decayed;
}
\end{verbatim}
For full credit, your code should follow the same general approach as the original code, rather than being a complete rewrite. 

\part[6]
Write a MIPS assembly function that implements your code from the previous part, but does not make any recursive calls. 
\end{parts}


\question[8]
\textbf{Days of the Week.}
Write a MIPS assembly function that takes as input a pointer to a list of 10 characters and an enum representing the days of the week.
The function should set the list to contain the word corresponding to that day of the week following C string guidelines. 
You might find it helpful to use the ASCII character representation to set the values. 
Assume that the enum is numbered from 0 to 6, starting on Sunday and ending on Saturday. 
The function should return whether the value passed was a valid enum value. 
For full credit, your function should use a jump table. 

\question[15]
\textbf{Overflow in a Tree.}
\begin{parts}
\part[8]
Write a MIPS assembly function that takes as input two double precision floating point numbers and returns a boolean (0 or 1) indicating whether they overflow (result in a special encoding number) when added together. 
\part[7]
Write a MIPS assembly function that takes as input a double precision floating point number and a pointer to a binary search tree node and returns a boolean indicating whether any value in the tree rooted at that node overflows when added to the floating point number. 
The tree nodes contains a double precision floating point number, a pointer to the left child node of the tree, and a pointer to the right child node of the tree, stored in that order. 
You can assume that any invalid pointer has the value zero. 
For full credit, your function should search the minimum number of tree nodes. 
For full credit, use your function from the previous part. 
\end{parts}

\question[10]
\textbf{Assembly Representation.}
Consider the following code:
\begin{verbatim}
    Loop:   sll $t0, $s0, 2 	
            addu $t1, $t0, $s3 	
            lw $t2, 0($t1) 
            blt $t2, $s2, Exit 
            addu $s0, $s0, $s1 
            j Loop 
    Exit:
\end{verbatim}
Convert the code from assembly language to machine code. 
Assume that the code starts at address 0x3080. 

\question[10]
\textbf{Writeup.}
In addition to your code, you should submit a brief (~1-2 pages) writeup that describes the state of your code. 
You should discuss what is working, what is not, and any known errors. 
In this discussion, you should reference the tests you did and how they support your claims. 
You will need to provide a brief description of your tests for you to reference. 
Your writeup should be in a separate .txt or .pdf file. 

\end{questions}

\end{document}